% Independently reviewed complete composite; two production assignments retained.
\paragraph{A finite ellipsoidal response at matched initial charge and energy.}
We additionally evolved finite, axisymmetric ellipsoidal preparations at
$\varepsilon=\pm0.02,\pm0.04$, with initial charge and full discrete energy
matched to each discretization's radial reference within the declared root
tolerance. Unlike the preceding second variation, this tests finite initial
data rather than energy compensation only through order $\varepsilon^2$.
At the fixed endpoint $T=32$ define
\[
 D(e)=\frac{C(+e)+C(-e)-2C(0)}{2e^2},\qquad C=Q_{\mathrm{core}}/Q_0.
\]
This is a symmetric final core response, not initial-subtracted net transport.
On the spatially refined grid,
$D(0.02)=0.0727746617$ and $D(0.04)=0.0730141978$, against the new
same-grid, same-timestep variational reference $c_{\mathrm{ref}}=0.0726948531$.
The prescribed heuristic envelope at $e=0.02$ is $1.2055721\times10^{-4}$;
the discrepancy $7.9808621\times10^{-5}$ plus this allowance meets the
predeclared one-percent agreement criterion. This envelope is not a certified
error bound or a confidence interval. The resolved-remainder criterion is
not met, so no measured Taylor order is claimed.

The separate core-fraction differences at $\varepsilon=+0.02,-0.02$ are
$2.8253326\times10^{-5}$ and $2.9966404\times10^{-5}$, respectively,
both positive in the stored endpoint data; the symmetric envelope is not
a separate certified bound for each sign. These are dimensionless charge-fraction
differences, not energy percentages; a positive symmetric coefficient alone
would not establish both signs. Figure~\ref{fig:fa-response} separates them.
The preparation changes geometry and radial chirp together. This finite-time
comparison establishes neither persistent enhancement, stationary formation,
binding nor general three-dimensional stability. The even-$\ell$, $m=0$
evolution excludes non-axisymmetric dynamics.

The complete twenty-trajectory matrix and separate variational reference
were assembled from two bounded production assignments: fourteen completed
trajectories and the reference were inherited unchanged, and six remaining
trajectories were completed separately. The first assignment retains its
\texttt{INCOMPLETE} status; the full matrix, not a favourable subset,
is used in the comparison.

\begin{figure}[htbp]
\centering
\includegraphics[width=.98\linewidth]{fa-response.pdf}
\caption{Finite symmetric endpoint coefficients for four discretizations
(left) and separate signed core-fraction differences on the spatially refined
grid (right). Caps are heuristic sensitivities, not confidence intervals.
Straight segments join stored time samples; no extrapolation or Taylor-order
fit is used. Initial redistribution is retained rather than subtracted.}
\label{fig:fa-response}
\end{figure}

\begin{figure}[htbp]
\centering
\includegraphics[width=.95\linewidth]{fa-fields-t0.pdf}
\caption{Stored initial amplitude (upper row) and relative phase (lower row)
for the radial and $\varepsilon=\pm0.04$ preparations. The meridional cut
uses $r\le12$ inside the computational domain $R=80$; the circle marks the
measurement radius $r=6$. The common amplitude scale is shared with
Fig.~\ref{fig:fa-fields-final}. Relative phase is
$\arg(\phi_\varepsilon\overline{\phi_0})$ in the common initial phase
convention, masked if either amplitude is below $10^{-4}$ of the shared
maximum. The centre $r<0.1$ is masked; grey is not zero amplitude or phase.}
\label{fig:fa-fields-initial}
\end{figure}

\begin{figure}[htbp]
\centering
\includegraphics[width=.95\linewidth]{fa-fields-t32.pdf}
\caption{The same stored-field reconstruction at $T=32$, with the common
scales and masks of Fig.~\ref{fig:fa-fields-initial}. Complex radial mode
coefficients are interpolated before evaluating the angular basis and phase.
These are axisymmetric amplitude/relative-phase cuts, not a general 3D
simulation, an absolute-phase observable, or a binding diagnostic.
Visually similar colours on the full $[-\pi,\pi]$ scale do not establish
exact phase equality.}
\label{fig:fa-fields-final}
\end{figure}

\paragraph{Phase rotation and the fixed-charge kinetic gap.}
A similar phase colour does not test rigid harmonic motion. We evaluated
necessary conditions on the six stored states with $\varepsilon=0,\pm0.04$
at $t=0,32$, using the complete $R=80$ box. For the reduced mode arrays
$U=r\phi$, $V=\partial_t U$, write
$\|X\|_h^2=4\pi h\sum_{j=1}^{J-1}\sum_\ell|X_{j\ell}|^2$,
$I=\|U\|_h^2$, $K=\|V\|_h^2$ and, for $I>0$,
$\omega_{\mathrm{fit}}=Q/(2I)$.
With $A_h(U)$ the actual discrete nonlinear acceleration, define
\[
 r_v=\frac{\|V-\ii\omega_{\mathrm{fit}}U\|_h^2}{K+\omega_{\mathrm{fit}}^2 I},
 \qquad
 r_a=\frac{\|A_h(U)+\omega_{\mathrm{fit}}^2 U\|_h^2}
                 {\|A_h(U)\|_h^2+\omega_{\mathrm{fit}}^4 I}.
\]
Rigid single-frequency motion requires both numerators to vanish. Initially
$r_v$ vanishes by construction, whereas $r_a\simeq0.0830$; the prescribed
phase velocity does not solve the stationary-profile equation. At $T=32$
the three states give $r_v=0.0282474$--$0.0282631$ and
$r_a=0.0586355$--$0.0588005$. These are squared-residual quotients, not
amplitude percentages or kinetic-energy fractions. This completed endpoint
postprocessing involved no further time integration. It provides neither a core-only classification,
a relaxation law, a certified continuum residual bound nor a general 3D
stability result.

The unnormalized velocity residual has a separate exact algebraic meaning.
In the positive canonical kinetic norm, minimizing velocities at fixed
instantaneous $U$ and total $Q$ gives $V_Q=\ii Q U/(2I)$ and
\[
 E_h(U,V)-E_{Q,h}[U]
 =J_v=\|V-\ii\omega_{\mathrm{fit}}U\|_h^2
 =K-\frac{Q^2}{4I},\qquad
 E_{Q,h}[U]=W_h(U)+\frac{Q^2}{4I},
\]
where $W_h$ contains the radial and angular gradient plus potential energy.
This is a comparison at fixed profile, not a physical projection or measured
relaxation: it generally changes regional charge and need not preserve
other conserved quantities. It does not minimize the spatial profile,
show how energy is radiated, or establish binding. In particular, $r_v$
normalizes $J_v$ by $K+Q^2/(4I)$, not by kinetic or total energy.
